% ═════════════════════════════════════════════════════════════════════════════
\section{Abstraction and Network Invariants}
% ═════════════════════════════════════════════════════════════════════════════

Abstraction-based methods tackle the parameterized problem by reducing it to a
finite, non-parameterized one. The idea is to construct an abstract system ---
typically a small, fixed-size system, that \emph{over-approximates} the
behavior of the full parameterized family. If a safety property holds in the
abstract system, it holds in all concrete instances~\cite{WL89,PXZ02}.

A particularly influential instantiation of this idea is the \emph{network
invariant} method~\cite{WL89}, further developed in~\cite{PXZ02}. Consider
a system consisting of $n$ copies of a process template running in parallel.
A network invariant is a single process~$I$ (or small composition of processes)
that captures the collective behavior of all $n-1$ processes that a given
process interacts with. Verifying the property then reduces to checking the
two-process composition of one concrete process with the invariant~$I$,
regardless of~$n$. If this finite system satisfies the property, and~$I$ is
a valid over-approximation of the environment, then all $n$-process instances
satisfy it too.

Finding the right invariant automatically is the central challenge. Approaches
range from manual construction guided by domain knowledge, to automated
inference using abstract interpretation, to counterexample-guided abstraction
refinement (CEGAR) adapted to the parameterized setting. In the CEGAR approach,
one begins with a coarse invariant, checks the finite abstract system, and if
a spurious counterexample is found, refines the invariant to rule it out. The
work of~\cite{PXZ02} combines this idea with \emph{$(0,1,\infty)$-counter
abstraction}, which tracks for each local state whether zero, one, or many
processes occupy it. This finite abstraction is often coarse but captures
enough structure to verify mutual exclusion and other coordination properties
automatically, and the abstraction can be systematically refined when needed.

The network invariant approach is closely related to view abstraction
(Section~3) in spirit: both methods replace the unbounded environment of a
process by a finite abstract summary. The key difference is in how that
summary is constructed and updated. Network invariants are defined globally
as a fixed process before verification begins; view abstraction constructs the
summary on the fly from the concrete reachability analysis, refining it
incrementally as the precision parameter~$k$ increases.
